\chapter{Background} \label{chap:background}


\section{Summarization Tasks}\label{sec:background-tasks}

Automatic generation of summaries is a well established task in the field of Natural Language Processing (NLP) that has been considered in a number of forms. Much early work in summarization focused on the news domain; multiple popular datasets have been introduced for the task of news summarization, such as the CNN/DailyMail~\cite{see-etal-2017-get} and XSUM~\cite{xsum-narayan-etal-2018-dont} datasets. Early work in summarization primarily considered \textit{summarization as information-compression}; the goal was to convey the same or most important information in a shorter amount of text~\cite{knight2002summarization, clarke2008global}. Summarization of scientific documents has taken the form of abstract generation~\cite{cohan2018discourse} and using additional context to augment abstracts, such as citations~\cite{Qazvinian2008ScientificPS, Cohan2015ScientificAS} or recordings of conference talks~\cite{Lev2019TalkSummAD}. In addition to the news and science domains, prior work has tackled specialized domains such as summarizing clinical notes~\cite{Abacha2023OverviewOT}, legal documents~\cite{Takale2023ASO}, or code summarization~\cite{Ahmad2020ATA}. Researchers have also considered multi-document summarization~\cite{DeYoung2021MS2MS, Fabbri2019MultiNewsAL}. Advances in language generation technology (see~\Cref{sec:background-methods}) have allowed researchers to consider more complex forms of summarization. More recent work has framed \textit{summarization as information-communication}, rather than information-compression. In this broader definition of summarization, we are able to consider a wider variety of tasks, goals, and users. For example, summarization can be used to answer queries~\cite{Nema2017DiversityDA, Fok2023QlarifyRE} or communicate complex topics to a general audience~\cite{goldsack2022making,Zaman2020HTSSAN}.  While researchers have imagined a large number of tasks that users may find useful in summary generation, each of these setups makes some assumptions about the user. They assume users simply want a shorter version of a news article or that all users will have similar background knowledge. These assumptions  constrain the user and may or may not apply to each person. In this thesis, my goal is to build and evaluate systems of summarization that prioritize the needs of the end-user.


\section{Summarization Methods}\label{sec:background-methods}

Methods in generating summaries can broadly be separated into two categories: extractive and abstractive. Extractive systems return text that is present in the input document as a summary while abstractive systems generate new text. Below I will provide an overview of each category.

\subsection{Extractive Summarization}

Methods are considered extractive if the system constructs a summary from the text in the input document. Many of these methods were born out of viewing summarization as information-compression, as discussed in~\Cref{sec:background-tasks}. For example, TextRank is an unsupervised method that extracts the most important sentences in a document based on the relative word count~\cite{Mihalcea2004TextRankBO}, which inspired PACSUM~\cite{Jie2023UnsupervisedES}, an extension of TextRank. In addition to unsupervised methods, prior work has considered supervised extractive methods. For example, BERTSumExt uses BERT as a sentence encoder augmented with inter-sentence Transformer layers to capture interactions between sentences~\cite{Liu2019TextSW}. The benefit of extractive methods is that they are simple by nature; they use the text within a document to return a summary. This means that extractive methods are more likely to produce factually correct summaries. Extractive summarization systems are also more interpretable; you can trace each sentence in the summary back to the original document. However, extractive methods are also more limiting, as they do not allow for transformations in the text. For example, extractive systems cannot produce a summary that is intended for a different audience than the original document. Modern systems may use extractive summarization methods as a lightweight, pre-processing step to shorten documents to fit in a context window~\cite{xu2020coarse, liu2023lost} or for other related tasks, such as document highlighting~\cite{Fok2023ScimIS}.

\subsection{Abstractive Summarization}

As extractive systems are limited in the types of outputs they can produce, abstractive summarization is a popular class of methods. Abstractive methods are able to generate new text that is not necessarily present in the original input document. This is more similar to how humans write summaries; human written summaries rarely copy sentences word for word. Additionally, abstractive summarization methods can transform the original text, allowing for more types of summaries. However, abstractive methods are also prone to hallucinations or factual inconsistencies~\cite{Ji2022SurveyOH}. In this thesis, I focus on developing abstractive summarization methods as abstraction is required for personalization; you cannot build a good summary for a non-expert only using the language present in a scientific paper.

Within the family of abstractive summarization, researchers have explored a number of methods. Prior work has experimented with Recurrent Neural Networks (RNNs)~\cite{Rafi2021RNNEA} and Pointer-Generator Networks~\cite{see-etal-2017-get}. First applied to the task of Machine Translation, the attention mechanism was first introduced by~\textcite{Bahdanau2014NeuralMT}; the core idea is that the model learns which words are important, or which words to ``attend'' to. \textcite{Vaswani2017AttentionIA} leveraged the attention mechanism in transformer networks, a type of neural network that uses self-attention to encode text. This innovation plus other innovations such as reinforcement-learning with human feedback (RLHF)~\cite{NIPS2017_d5e2c0ad, NEURIPS2022_b1efde53}, instruction tuning~\cite{wei2022finetuned,chung2022scaling}, and even simple scaling of models~\cite{kaplan2020scaling, hoffmann2022training, Wei2022EmergentAO}, have led to a sharp increase in the overall performance of language models, resulting in the modern era of language technologies. 

Language models can broadly be categorized into three types: encoder-only, decoder-only, or encoder-decoder. Encoder-only models do not generate text, so they are not applicable to the task of abstractive summarization; these models produce representations of language, which can be then used in downstream tasks such as text classification. Decoder-only models use a language modeling objective to generate text. In other words, given the previous tokens, generate the next token. In the case of summarization, the input document is provided to the LM which is then prompted to generate a summary~\cite{Zhang2023BenchmarkingLL}. Popular decoder-only language models such as GPT4, Mistral, and Llama have been used for summarization~\cite{Zhang2023BenchmarkingLL,cachola-etal-2024-knowledge,Laskar2023BuildingRM,Laskar2023ASS}. Encoder-decoder models first encode the input text (the input document) to create a representation of that text. The intermediate representation is then passed to a decoder to generate the output text (summary). Popular encoder-decoder models include BART~\cite{lewis2020bart} and T5~\cite{Raffel2019ExploringTL}. In general, at the same size, encoder-decoder language models tend to perform better on summarization tasks. However, larger decoder-only models have shown impressive performance on summarization benchmarks~\cite{goyal2022news,Zhang2023BenchmarkingLL}. Since there are no encoder-decoder models that compare in size to the largest decoder-only models, it is unclear how a larger encoder-decoder model would compare in performance. Decoder-only models are generally more popular, as they are better able to leverage unstructured data during pre-training and are generalizable to more tasks.
In this thesis, I generally focus on decoder-only language models for their flexibility, although I make an effort to focus on smaller, lower compute models when possible to maintain accessibility of my methodology. 

\section{Evaluation of Summaries}\label{sec:background-eval-summaries}
Given the natural variance in human written summaries, automatic evaluation of summaries is a challenging problem. There are a number of factors that influence the quality of a summary, such as information content, fluency, factual correctness, etc. Considering these challenges, human evaluation remains the gold standard; if we want to know if a summary is acceptable to humans, we should ask them. However, human evaluation is too slow and expensive to use for iterating on methodology, so it is still pertinent that we design better automatic evaluation metrics on which to base our modeling decisions. Below, I will describe the current state of summarization evaluation. 

\subsection{Quality Evaluation}
The majority of metrics for summarization try to measure the overall quality of the summary. These metrics broadly fall into two categories: reference-based and reference-free.  Reference-based metrics use a ground truth, generally a human written summary, as a reference to compare with the generated summary. The goal of reference-based metrics is to compare how similar the generated summary is to the reference summary. Reference-based metrics can measure either lexical similarity or semantic similarity. Popular lexical similarity metrics include ROUGE~\cite{lin2004rouge}, METEOR~\cite{banerjee2005meteor}, and BLEU~\cite{papineni2002bleu}. These metrics are popular given their simplicity and high interpretability; it is easy to understand why your system received a high or low score. However, the simplicity of these metrics is also their disadvantage. For example, the sentences ``Isabel bought a soda then went to the movies'' and ``Isabel purchased a Coke before walking to the theater'' have very similar semantic meanings but very little lexical overlap. To address this concern, prior work has also introduced evaluation metrics that aim to measure the semantic similarity of a generated summary to the reference summary. Popular semantic similarity metrics include MoverScore~\cite{zhao2019moverscore}, BERTScore~\cite{zhang2020bertscore}, and BARTScore~\cite{Yuan2021BARTScoreEG}. While these metrics tend to correlate more with human judgments than lexical similarity metrics, they also are not as interpretable and are subject to bias~\cite{Sun2022BERTScoreIU}.

Reference based metrics require a reference summary. This poses two issues: (1) reference summaries can be expensive to collect and (2) it assumes that the reference summary is the best summary, which may or may not be a correct assumption. Reference free evaluation metrics try to address these issues by providing a measure of quality without a reference, instead measuring how well the summary captures the input document. For example, BLANC scores use Masked Language Models to measure how well the summary can be used to fill in the blanks of the original document~\cite{vasilyev2020blanc}. Other researchers have considered prompting LLMs to rate the quality of a summary~\cite{liu2023geval}. While these metrics have shown relatively high correlation with human judgment, they similarly are uninterpretable and subject to bias~\cite{he-etal-2023-blind, Shen2023LargeLM, Liu2023LLMsAN}.

In general, there is no perfect quality measure of a summary. It is best practice to use a combination of different metrics to evaluate the quality of a system. In this thesis, I use a number of these general quality metrics as baselines, as they are the standard in summarization evaluation. However, they fall short of the vision of user-centric evaluation. A generally high quality summary might still be unusable to a specific person if that summary assumes the wrong background or is targeted at a different use case. 

\subsection{Specific attributes}
In addition to overall quality, there are other attributes we may be interested in measuring to assess the performance of a system. For example, Unieval provides scores for the coherence, consistency, fluency, and relevance of a summary, in addition to an overall quality score~\cite{Zhong2022TowardsAU}. For the task of lay summarization, we are interested in the readability of a summary, in addition to the overall quality. Popular readability metrics include the Flesch Kincaid Grade Level (FKGL) and Dale Chall Readability Score (DCRS)~\cite{goldsack2022making}. FKGL provides a score that is a combination of the number of words in the sentences and the number of syllables in the words. DCRS uses a dictionary of common words and returns a score based on the frequency of common words. In addition to token-based metrics, researchers have also considered training metrics to predict readability~\cite{Imperial2021BERTEF}.
FActScore is a metric that predicts the factual correctness of generated text~\cite{min2023factscore}. These metrics consider aspects of summarization that are important to users other than overall quality, allowing researchers to discern where the systems have room for improvement along a more specific axis than general quality. Throughout this thesis, I use a number of these metrics to measure the efficacy of my methods along specific axes of interest. I also build-off of past work in summarization evaluation to introduce methods of evaluation that account for user-specific interests.